\chapter{Methodology}

This chapter presents the proposed hybrid sparse-sampling architecture for reasoning video segmentation.
The framework operates in three stages: (1) keyframe selection based on inter-frame motion analysis,
(2) semantic segmentation of selected keyframes using a multimodal large language model, and
(3) mask reconstruction for non-keyframe frames via Universal Kriging interpolation on contour-vectorized
representations.

\section{System Overview}

Given a video sequence $V = \{f_1, f_2, \ldots, f_N\}$ of $N$ frames and a natural language query $q$,
the goal is to produce a binary segmentation mask $m_t$ for each frame $t \in \{1, \ldots, N\}$ that
corresponds to the object described by $q$. Rather than applying the semantic model to every frame,
our method selects a sparse subset of keyframes $K \subset \{1, \ldots, N\}$ with $|K| \ll N$,
processes only these frames through the expensive semantic model, and reconstructs the remaining
masks through geostatistical interpolation.

The overall pipeline is:
\begin{enumerate}
    \item \textbf{Keyframe Selection:} Compute inter-frame motion magnitudes and select a subset
          $K$ of keyframes at a target sampling rate $r$ (e.g., $r = 0.2$).
    \item \textbf{Semantic Segmentation:} Apply the reasoning segmentation model (e.g., VideoLISA)
          to keyframes $\{f_k : k \in K\}$ to obtain ground-truth-quality masks $\{m_k : k \in K\}$.
    \item \textbf{Contour Vectorization:} Convert each binary keyframe mask $m_k$ to a fixed-length
          polygon representation $\mathbf{v}_k \in \mathbb{R}^{P \times 2}$ with $P$ vertices.
    \item \textbf{Kriging Interpolation:} For each vertex $i \in \{1, \ldots, P\}$, perform
          1D Ordinary Kriging along the time axis to predict vertex positions at non-keyframe times.
    \item \textbf{Rasterization:} Convert interpolated polygon vertices back to binary masks.
\end{enumerate}

\section{Keyframe Selection}
\label{sec:keyframe-selection}

The keyframe selector identifies frames where significant visual change occurs, ensuring that
the semantic model is applied at temporally informative moments.

\subsection{Motion Estimation}

For each consecutive frame pair $(f_t, f_{t+1})$, we compute a scalar motion magnitude $\phi_t$.
Three motion estimation backends are supported:

\begin{itemize}
    \item \textbf{Farneback optical flow}: Dense optical flow via polynomial expansion~\cite{farneback_2003}.
          The magnitude is the mean pixel-wise flow norm:
          $\phi_t = \frac{1}{HW} \sum_{x,y} \|\mathbf{u}_{t}(x,y)\|_2$
          where $\mathbf{u}_t$ is the flow field between frames $t$ and $t+1$.
    \item \textbf{DIS optical flow}: Dense Inverse Search~\cite{kroeger_dis_2016}, offering faster
          computation at comparable accuracy.
    \item \textbf{Pixel difference}: Mean absolute grayscale difference,
          $\phi_t = \frac{1}{HW} \sum_{x,y} |g_t(x,y) - g_{t+1}(x,y)|$,
          where $g_t$ is the grayscale conversion of $f_t$. This serves as a GPU-free fallback.
\end{itemize}

\subsection{Selection Strategies}

Given the motion magnitudes $\{\phi_1, \ldots, \phi_{N-1}\}$ and a target sampling rate $r$,
the selector chooses $k = \lceil r \cdot N \rceil$ keyframes. We propose three strategies:

\paragraph{Top-K (baseline).} Select the $k$ frames with the highest preceding motion magnitude.
This approach clusters keyframes around motion bursts, leaving long static stretches without
temporal support for interpolation.

\paragraph{Segment-based (proposed).} Divide the video into $k - 2$ equal temporal segments
(excluding the always-included first and last frames). Within each segment, select the frame
with the highest motion magnitude. This ensures even temporal coverage while still capturing
local motion peaks within each segment.

\paragraph{Local maxima.} Identify local maxima in the motion signal $\{\phi_t\}$ and greedily
select the strongest peaks subject to a minimum gap constraint
$\Delta_{\min} = \lfloor N / (k + 1) \rfloor$ between any two selected keyframes. This prevents
clustering while preferring perceptually meaningful transition points.

In all strategies, the first frame ($t = 1$) and last frame ($t = N$) are always included
to provide boundary conditions for interpolation.

\section{Contour Vectorization}
\label{sec:contour-vectorization}

Binary segmentation masks must be converted to a representation amenable to continuous
interpolation. We adopt a contour-based approach: each mask is represented as a fixed-length
polygon with $P$ vertices (default $P = 64$).

\subsection{Contour Extraction and Resampling}

Given a binary mask $m_k$, we extract the largest connected contour using boundary tracing.
The raw contour is then resampled to exactly $P$ vertices by linear interpolation along the
arc-length parameterization:
\begin{equation}
    s_j = \frac{j}{P} \cdot L, \quad j = 0, 1, \ldots, P-1
\end{equation}
where $L$ is the total perimeter of the contour. The resampled vertex at arc-length $s_j$ is
obtained by linear interpolation between the two nearest original contour points.

\subsection{Centroid-Aligned Vertex Ordering}

To ensure consistent vertex correspondence across keyframes, we apply centroid-aligned ordering.
After resampling, the vertex list is cyclically rotated so that vertex 0 is the point closest
to the 12 o'clock position relative to the contour centroid (maximum $-\Delta y / \|\Delta\|$
direction). This prevents rotational drift in the vertex indexing across frames, ensuring that
vertex $i$ at keyframe $k_1$ corresponds to approximately the same physical location as vertex
$i$ at keyframe $k_2$.

\section{Universal Kriging Interpolation}
\label{sec:kriging-interpolation}

Kriging is a geostatistical interpolation method that provides the Best Linear Unbiased
Prediction (BLUP) by exploiting spatial correlation structure through variogram modeling.
We apply it to temporal interpolation of contour vertices.

\subsection{Per-Vertex 1D Kriging}

For each vertex index $i \in \{1, \ldots, P\}$, we treat the $x$-coordinate time series
$\{v^{(i)}_x(k) : k \in K\}$ and $y$-coordinate time series $\{v^{(i)}_y(k) : k \in K\}$
as independent 1D spatial processes along the time axis. For target frame $t \notin K$, the
Kriging predictor is:
\begin{equation}
    \hat{v}^{(i)}_x(t) = \sum_{j=1}^{|K|} \lambda_j(t) \cdot v^{(i)}_x(k_j)
\end{equation}
where the weights $\lambda_j(t)$ are determined by solving the Kriging system:
\begin{equation}
    \begin{bmatrix} \gamma(k_1, k_1) & \cdots & \gamma(k_1, k_{|K|}) & 1 \\
                    \vdots & \ddots & \vdots & \vdots \\
                    \gamma(k_{|K|}, k_1) & \cdots & \gamma(k_{|K|}, k_{|K|}) & 1 \\
                    1 & \cdots & 1 & 0 \end{bmatrix}
    \begin{bmatrix} \lambda_1 \\ \vdots \\ \lambda_{|K|} \\ \mu \end{bmatrix} =
    \begin{bmatrix} \gamma(k_1, t) \\ \vdots \\ \gamma(k_{|K|}, t) \\ 1 \end{bmatrix}
\end{equation}
where $\gamma(h)$ is the variogram function and $\mu$ is the Lagrange multiplier enforcing
unbiasedness.

\subsection{Variogram Models}

The variogram $\gamma(h) = \frac{1}{2} \text{Var}[Z(t+h) - Z(t)]$ describes the spatial
correlation structure. We evaluate four standard models:

\begin{itemize}
    \item \textbf{Gaussian:} $\gamma(h) = c \left[1 - \exp\left(-\frac{h^2}{a^2}\right)\right]$
    \item \textbf{Exponential:} $\gamma(h) = c \left[1 - \exp\left(-\frac{h}{a}\right)\right]$
    \item \textbf{Spherical:} $\gamma(h) = c \left[\frac{3h}{2a} - \frac{h^3}{2a^3}\right]$ for $h \leq a$, $c$ otherwise
    \item \textbf{Linear:} $\gamma(h) = c \cdot \frac{h}{a}$ for $h \leq a$, $c$ otherwise
\end{itemize}
where $c$ is the sill (total variance) and $a$ is the range parameter. The variogram parameters
are fitted automatically from the observed keyframe data. Based on our ablation study
(Section~\ref{sec:variogram-ablation}), the exponential model is used as the default.

\subsection{Rasterization}

After Kriging interpolation, the predicted vertices $\hat{\mathbf{v}}_t \in \mathbb{R}^{P \times 2}$
for each non-keyframe $t$ are rasterized back to a binary mask using polygon fill operations.
Keyframe masks are passed through unchanged, preserving the original semantic model output at
those frames.
